%==============================================================
%  lecons/analyse/jacobi.tex — Les polynômes de Jacobi
%==============================================================

\lecontitle{Les polynômes de Jacobi}{Analyse réelle — Polynômes orthogonaux et approximation}

%=== NOTATION LOCALE =========================================
\newcommand{\Jac}{P_n^{(\alpha,\beta)}}
\newcommand{\wab}{w_{\alpha,\beta}}

%=============================================================
%  § 1 — DÉFINITION ET RÉSULTATS FONDAMENTAUX
%=============================================================
\secnum{navyblue}{1}{Définition et résultats fondamentaux}

\begin{tcolorbox}[thmbox]
  \textbf{Définition.}\enspace\index{polynômes de Jacobi}
  Soient $\alpha,\beta>-1$. Sur $\mathbb{R}[X]$, on munit l'intervalle
  $[-1,1]$ du poids
  \[
    \wab(x) \;=\; {(1-x)}^{\alpha}\,{(1+x)}^{\beta},
  \]
  intégrable sur $]-1,1[$ précisément lorsque $\alpha,\beta>-1$. Le
  \emph{$n$-ième polynôme de Jacobi} $\Jac$ est défini par la
  \emph{formule de Rodrigues}\index{formule de Rodrigues} généralisée~:
  \[
    \Jac(x) \;=\; \frac{{(-1)}^n}{2^n\,n!}\,
    {(1-x)}^{-\alpha}{(1+x)}^{-\beta}\,
    \frac{\mathrm{d}^n}{\mathrm{d}x^n}\!
    \Bigl[{(1-x)}^{n+\alpha}{(1+x)}^{n+\beta}\Bigr].
  \]
  C'est un polynôme de degré $n$, normalisé par
  $\Jac(1)=\dbinom{n+\alpha}{n}$, où
  $\dbinom{n+\alpha}{n}=\dfrac{\Gamma(n+\alpha+1)}{\Gamma(\alpha+1)\,n!}$
  désigne le coefficient binomial généralisé ($\alpha$ réel).

  \medskip
  \textbf{Théorème (propriétés fondamentales).}\enspace%
  \index{polynômes orthogonaux}
  La famille ${\bigl(\Jac\bigr)}_{n\geq 0}$ vérifie~:
  \begin{itemize}
    \item \textit{Orthogonalité} sur $[-1,1]$ pour le poids $\wab$~:
    \[
      \int_{-1}^{1} P_m^{(\alpha,\beta)}(x)\,P_n^{(\alpha,\beta)}(x)\,
      \wab(x)\,\mathrm{d}x = 0 \qquad (m\neq n),
    \]
    et, pour $m=n$, la norme vaut
    \[
      h_n^{(\alpha,\beta)}=\frac{2^{\alpha+\beta+1}}{2n+\alpha+\beta+1}\,
      \frac{\Gamma(n+\alpha+1)\,\Gamma(n+\beta+1)}
           {n!\,\Gamma(n+\alpha+\beta+1)}.
    \]
    \item \textit{Équation différentielle de Jacobi}~: $y=\Jac$ est
          solution de
    \[
      (1-x^2)\,y'' + \bigl[\beta-\alpha-(\alpha+\beta+2)\,x\bigr]\,y'
      + n(n+\alpha+\beta+1)\,y = 0.
    \]
    \item \textit{Récurrence à trois termes}~: il existe des suites de
          réels $(a_n),(b_n),(c_n),(d_n)$ telles que
    \[
      a_n\,P_{n+1}^{(\alpha,\beta)}(x)
      = (b_n+x)\,c_n\,P_n^{(\alpha,\beta)}(x)
      - d_n\,P_{n-1}^{(\alpha,\beta)}(x),
    \]
    avec $P_0^{(\alpha,\beta)}=1$ et
    $P_1^{(\alpha,\beta)}=\tfrac12(\alpha+\beta+2)x+\tfrac12(\alpha-\beta)$.
    \item $\Jac$ possède $n$ racines simples dans $\,]-1,1[\,$ (nœuds de
          la quadrature de Gauss associée au poids $\wab$).
  \end{itemize}

  \medskip
  \textbf{Famille mère.}\enspace Les familles classiques sur $[-1,1]$
  s'obtiennent par spécialisation du couple $(\alpha,\beta)$~:
  \[
    \begin{array}{lll}
      \text{Legendre } (P_n): & \alpha=\beta=0, & w\equiv 1,\\[2pt]
      \text{Tchebychev } 1^{\text{re}}\text{ esp. } (T_n):
        & \alpha=\beta=-\tfrac12, & {(1-x^2)}^{-1/2},\\[2pt]
      \text{Tchebychev } 2^{\text{e}}\text{ esp. } (U_n):
        & \alpha=\beta=+\tfrac12, & {(1-x^2)}^{+1/2},\\[2pt]
      \text{Gegenbauer } (C_n^{\lambda}):
        & \alpha=\beta=\lambda-\tfrac12, & {(1-x^2)}^{\lambda-1/2}.
    \end{array}
  \]
  Les polynômes de Gegenbauer\index{polynômes de Gegenbauer} (ou
  \emph{ultrasphériques}\index{polynômes ultrasphériques}), cas symétrique
  $\alpha=\beta$, contiennent eux-mêmes Legendre ($\lambda=\tfrac12$) et les
  deux familles de Tchebychev (à normalisation près).
\end{tcolorbox}

\begin{significationintuitive}
Un \emph{seul} cadre — l'orthogonalité sur $[-1,1]$ pour le poids
$\wab(x)={(1-x)}^\alpha{(1+x)}^\beta$ — engendre toutes les familles
classiques de cet intervalle. Les exposants $\alpha$ et $\beta$ règlent le
comportement du poids \emph{aux bords} $x=1$ et $x=-1$~: positif, il
\emph{déprime} la fonction (les racines fuient les bords), négatif (mais
$>-1$), il \emph{renforce} le poids et concentre l'information près des
extrémités. Faire varier $(\alpha,\beta)$, c'est donc choisir où l'on
souhaite le mieux approcher~: poids uniforme pour Legendre, poids
singulier aux bords pour Tchebychev. Les polynômes de Jacobi sont ainsi
le « tronc commun » dont Legendre, Tchebychev et Gegenbauer ne sont que
des branches.
\end{significationintuitive}

\decsep{}

%=============================================================
%  § 2 — DÉMONSTRATION
%=============================================================
\secnum{navyblue}{2}{Démonstration}

\begin{ideepreuve}
L'orthogonalité se prouve comme pour Legendre, en intégrant $n$ fois par
parties la formule de Rodrigues~: les termes de bord s'annulent car les
exposants $n+\alpha$ et $n+\beta$ sont $>0$, ce qui force la nullité aux
points $\pm 1$. La spécialisation aux familles connues se lit directement
sur le poids $\wab$ en substituant les valeurs de $(\alpha,\beta)$.
L'équation différentielle découle de Rodrigues, et la structure générale
(récurrence à trois termes, racines simples réelles) est une conséquence
abstraite de l'orthogonalité, valable pour toute famille à poids positif.
\end{ideepreuve}

\vspace{10pt}
\rubriq{Preuve}

\begin{mdframed}[style=proofbar]
  Posons $\rho_n(x)={(1-x)}^{n+\alpha}{(1+x)}^{n+\beta}$, de sorte que
  $\Jac=\dfrac{{(-1)}^n}{2^n n!}\,\wab^{-1}\rho_n^{(n)}$. Comme
  $n+\alpha>0$ et $n+\beta>0$, le facteur $\rho_n$ et ses dérivées
  d'ordre $\leq n-1$ s'annulent en $\pm 1$~:
  $\rho_n^{(k)}(\pm 1)=0$ pour $0\leq k\leq n-1$.

  \medskip
  \textit{Étape 1 — Orthogonalité (Rodrigues + IPP).}
  Soit $Q$ un polynôme de degré $m<n$. Alors
  \[
    \int_{-1}^{1}\Jac(x)\,Q(x)\,\wab(x)\,\mathrm{d}x
    = \frac{{(-1)}^n}{2^n n!}\int_{-1}^{1}\rho_n^{(n)}(x)\,Q(x)\,\mathrm{d}x.
  \]
  En intégrant $n$ fois par parties, les crochets
  $\bigl[\rho_n^{(k)}\,Q^{(n-1-k)}\bigr]_{-1}^{1}$ sont nuls (car
  $\rho_n^{(k)}(\pm 1)=0$ pour $k\leq n-1$), d'où
  \[
    \int_{-1}^{1}\rho_n^{(n)}\,Q\,\mathrm{d}x
    = {(-1)}^n\int_{-1}^{1}\rho_n\,Q^{(n)}\,\mathrm{d}x = 0,
  \]
  puisque $Q^{(n)}\equiv 0$. Ainsi $\Jac$ est orthogonal (pour le poids
  $\wab$) à tout polynôme de degré $<n$ ; en particulier
  $\langle P_m^{(\alpha,\beta)},P_n^{(\alpha,\beta)}\rangle_{\wab}=0$ si
  $m\neq n$. Pour la norme, on prend $Q=\Jac$, dont le coefficient dominant
  vaut
  \[
    k_n=\frac{\Gamma(2n+\alpha+\beta+1)}{2^n\,n!\,\Gamma(n+\alpha+\beta+1)},
    \qquad\text{d'où}\qquad Q^{(n)}=n!\,k_n ;
  \]
  en explicitant la dernière intégrale
  $\int_{-1}^1\rho_n\,\mathrm{d}x
  =2^{2n+\alpha+\beta+1}\,B(n+\alpha+1,\,n+\beta+1)$ par la fonction Bêta,
  on obtient la norme $h_n^{(\alpha,\beta)}$ annoncée.

  \medskip
  \textit{Étape 2 — Spécialisation aux familles connues.}
  Le poids $\wab(x)={(1-x)}^\alpha{(1+x)}^\beta$ se simplifie selon le
  couple $(\alpha,\beta)$~:
  \[
    \alpha=\beta=0\ :\ \wab\equiv 1\ \text{(Legendre)} ; \quad
    \alpha=\beta=\pm\tfrac12\ :\ \wab={(1-x^2)}^{\pm 1/2}\ \text{(Tchebychev)}.
  \]
  Pour $\alpha=\beta=-\tfrac12$, le poids ${(1-x^2)}^{-1/2}$ est celui des
  $T_n$ ; le changement $x=\cos\theta$ donne $\mathrm{d}x=-\sin\theta\,
  \mathrm{d}\theta$ et $\wab\,\mathrm{d}x=-\mathrm{d}\theta$, d'où
  $\int_{-1}^1 T_mT_n\,\wab\,\mathrm{d}x=\int_0^\pi\cos(m\theta)\cos(n\theta)
  \,\mathrm{d}\theta$, nul pour $m\neq n$~: on retrouve $T_n(\cos\theta)
  =\cos(n\theta)$. Pour $\alpha=\beta=+\tfrac12$, le poids
  ${(1-x^2)}^{1/2}$ donne de même $U_n(\cos\theta)=\sin((n+1)\theta)/
  \sin\theta$. À chaque substitution, l'unicité (à un facteur près) du
  système orthogonal identifie $\Jac$ à la famille correspondante.

  \medskip
  \textit{Étape 3 — Équation différentielle.}
  L'opérateur $\mathcal{L}(y)=\bigl(\wab(1-x^2)\,y'\bigr)'$ est, par
  construction du poids, \emph{auto-adjoint} pour $\langle\cdot,\cdot
  \rangle_{\wab}$. On vérifie sur Rodrigues que $\mathcal{L}(\Jac)
  =-\,n(n+\alpha+\beta+1)\,\wab\,\Jac$, ce qui, après division par $\wab$
  et développement, donne exactement
  $(1-x^2)y''+[\beta-\alpha-(\alpha+\beta+2)x]y'+n(n+\alpha+\beta+1)y=0$.
  Ainsi $\Jac$ est vecteur propre de $\mathcal{L}$.

  \medskip
  \textit{Étape 4 — Structure générale (récurrence et racines).}
  Ces propriétés ne dépendent que de l'orthogonalité pour un poids positif,
  non du choix particulier de $\wab$.
  \emph{Récurrence à trois termes~:} $x\,\Jac$ est de degré $n+1$, donc se
  décompose sur la base $\bigl(P_k^{(\alpha,\beta)}\bigr)_{k\leq n+1}$ ; pour
  $k\leq n-2$, $\deg(xP_k^{(\alpha,\beta)})\leq n-1<n$ entraîne que le
  coefficient est nul. Il ne subsiste que les indices $n-1,n,n+1$, d'où une
  relation de la forme annoncée.
  \emph{Racines~:} $\Jac$ étant orthogonal à $1$ pour le poids \emph{positif}
  $\wab$, on a $\int_{-1}^1\Jac\,\wab=0$, donc $\Jac$ change de signe dans
  $]-1,1[$. Si $x_1,\dots,x_r$ ($r<n$) sont ses changements de signe, le
  polynôme $R=\prod(X-x_i)$ vérifie $\deg R<n$, donc
  $\langle\Jac,R\rangle_{\wab}=0$ ; mais $\Jac\,R\,\wab$ garde un signe
  constant et n'est pas identiquement nul~: contradiction. Donc $r=n$ et
  $\Jac$ a $n$ racines simples dans $]-1,1[$.
  \hfill$\blacksquare$
\end{mdframed}

\decsep{}

%=============================================================
%  § 3 — EXEMPLES, CONTRE-EXEMPLES, EXERCICES
%=============================================================
\secnum{exgreen}{3}{Exemples, contre-exemples et exercices}

\rubriq{Exemples}

\begin{mdframed}[style=exbar]
  \textbf{1. Tableau des spécialisations.}\enspace Selon le couple
  $(\alpha,\beta)$, $\Jac$ redonne (à normalisation près) une famille
  classique sur $[-1,1]$~:
  \[
    \begin{array}{llll}
      (\alpha,\beta) & \text{Famille} & \text{Poids } \wab(x) & \text{Normalisation}\\[2pt]
      (0,0)            & \text{Legendre } P_n & 1 & P_n(1)=1\\[2pt]
      (-\tfrac12,-\tfrac12) & \text{Tchebychev } T_n & {(1-x^2)}^{-1/2} & T_n(\cos\theta)=\cos n\theta\\[2pt]
      (+\tfrac12,+\tfrac12) & \text{Tchebychev } U_n & {(1-x^2)}^{+1/2} & U_n(\cos\theta)=\tfrac{\sin((n+1)\theta)}{\sin\theta}\\[2pt]
      (\lambda-\tfrac12,\lambda-\tfrac12) & \text{Gegenbauer } C_n^{\lambda} & {(1-x^2)}^{\lambda-1/2} & \text{ultrasphérique}
    \end{array}
  \]

  \smallskip\noindent
  \textbf{2. Tchebychev de $2^{\text{e}}$ espèce.}\enspace
  Pour $\alpha=\beta=\tfrac12$, le changement $x=\cos\theta$ donne
  \[
    U_n(\cos\theta)=\frac{\sin\bigl((n+1)\theta\bigr)}{\sin\theta},
  \]
  d'où $U_0=1$, $U_1=2x$, $U_2=4x^2-1$, $U_3=8x^3-4x$. Ces polynômes sont
  orthogonaux pour le poids ${(1-x^2)}^{1/2}$ et vérifient la récurrence
  $U_{n+1}=2xU_n-U_{n-1}$.

  \smallskip\noindent
  \textbf{3. Gegenbauer.}\enspace\index{polynômes de Gegenbauer}
  Le cas symétrique $\alpha=\beta=\lambda-\tfrac12$ définit les
  $C_n^{\lambda}$, de fonction génératrice
  \[
    {(1-2xt+t^2)}^{-\lambda}=\sum_{n=0}^{+\infty}C_n^{\lambda}(x)\,t^n
    \qquad (\lambda>0).
  \]
  Pour $\lambda=\tfrac12$ on retombe sur la génératrice des $P_n$ : les
  ultrasphériques\index{polynômes ultrasphériques} interpolent ainsi entre
  Tchebychev et Legendre.
\end{mdframed}

\vspace{6pt}
\rubriq{Contre-exemples et pièges}

\begin{mdframed}[style=cexbar]
  \textbf{Condition $\alpha,\beta>-1$.}\enspace
  Le poids $\wab(x)={(1-x)}^\alpha{(1+x)}^\beta$ n'est intégrable sur
  $]-1,1[$ que si $\alpha>-1$ \emph{et} $\beta>-1$~: pour $\alpha\leq -1$,
  $\int_{-1}^1{(1-x)}^\alpha\,\mathrm{d}x$ diverge au voisinage de $1$, et le
  produit scalaire perd tout sens. La famille de Jacobi n'est définie qu'à
  l'intérieur de ce domaine.

  \smallskip\noindent
  \textbf{Hermite et Laguerre ne sont pas des Jacobi.}\enspace
  Les polynômes d'Hermite (poids $e^{-x^2}$ sur $\mathbb{R}$) et de Laguerre
  (poids $e^{-x}$ sur $[0,+\infty[$) sont orthogonaux sur un intervalle
  \emph{non borné}~: ils n'entrent pas dans le cadre Jacobi, qui vit sur le
  segment $[-1,1]$. On les obtient seulement comme \emph{cas limites}
  (par exemple Hermite via $C_n^{\lambda}$ quand $\lambda\to+\infty$ après
  changement d'échelle), jamais par spécialisation directe de $(\alpha,\beta)$.

  \smallskip\noindent
  \textbf{Normalisations divergentes.}\enspace
  La convention $\Jac(1)=\binom{n+\alpha}{n}$ ne coïncide pas avec celles de
  $T_n$, $U_n$ ou $C_n^{\lambda}$~: passer d'une famille à l'autre demande un
  facteur multiplicatif. On a par exemple $P_n=P_n^{(0,0)}$ mais
  $T_n=\dfrac{n!}{(1/2)_n}\,P_n^{(-1/2,-1/2)}$, où
  $(a)_n=a(a+1)\cdots(a+n-1)=\dfrac{\Gamma(a+n)}{\Gamma(a)}$ désigne le
  \emph{symbole de Pochhammer}\index{symbole de Pochhammer}~: ne jamais
  identifier deux familles sans expliciter le facteur d'échelle.
\end{mdframed}

\vspace{6pt}
\exolabel

\begin{mdframed}[style=exobar]
  \begin{enumerate}
    \item \textit{(Legendre.)} En posant $\alpha=\beta=0$ dans la formule de
          Rodrigues de Jacobi, vérifier que l'on retrouve
          $P_n(x)=\dfrac{1}{2^n n!}\dfrac{\mathrm{d}^n}{\mathrm{d}x^n}
          {(x^2-1)}^n$, c'est-à-dire les polynômes de Legendre.

    \item \textit{(Tchebychev $1^{\text{re}}$ espèce.)} Pour
          $\alpha=\beta=-\tfrac12$, montrer via $x=\cos\theta$ que le poids
          ${(1-x^2)}^{-1/2}$ rend orthogonale la famille
          $T_n(\cos\theta)=\cos(n\theta)$, et vérifier la récurrence
          $T_{n+1}=2xT_n-T_{n-1}$.

    \item \textit{(Tchebychev $2^{\text{e}}$ espèce.)} Pour
          $\alpha=\beta=+\tfrac12$, établir
          $U_n(\cos\theta)=\sin((n+1)\theta)/\sin\theta$, calculer
          $U_2,U_3,U_4$, et vérifier
          $\int_{-1}^1 U_mU_n\,{(1-x^2)}^{1/2}\,\mathrm{d}x=\tfrac{\pi}{2}
          \delta_{mn}$.

    \item \textit{(Équation différentielle.)} Vérifier que pour
          $\alpha=\beta=0$ l'équation de Jacobi se réduit à l'équation de
          Legendre $(1-x^2)y''-2xy'+n(n+1)y=0$.

    \item \textit{(Gegenbauer.)} À partir de la fonction génératrice
          ${(1-2xt+t^2)}^{-\lambda}=\sum_n C_n^{\lambda}(x)t^n$, calculer
          $C_0^{\lambda},C_1^{\lambda},C_2^{\lambda}$ et retrouver les $P_n$
          pour $\lambda=\tfrac12$.

    \item \textit{(Intégrabilité.)} Montrer que
          $\int_{-1}^1{(1-x)}^\alpha{(1+x)}^\beta\,\mathrm{d}x
          =2^{\alpha+\beta+1}B(\alpha+1,\beta+1)$ pour $\alpha,\beta>-1$, et
          en déduire que la condition $\alpha,\beta>-1$ est nécessaire et
          suffisante à la convergence. \textit{Indication :} poser
          $x=2u-1$.
  \end{enumerate}
\end{mdframed}

\leconfooter{C.\ G.\ J.\ Jacobi, \textit{Untersuchungen über die Differentialgleichung der hypergeometrischen Reihe} (1859, posthume)}
